\section{22 августа 1914 г.}

\lp{20260604_194241882}

Милый мой, дорогой Миля! Мы с Колей так по вас стосковались, и нас так мучит
неизвестность один ли вы или с кем-нибудь, что решили дать вам телеграмму. Не
знаем даже доходит ли все это до вас. Дошла ли ответная наша телеграмма и
мамина телеграмма о папе, и все мои письма, и Колино, папашино,
Верочкино\textellipsis

Не нужны ли вам деньги? Теперь есть возможность переслать из России. Мы сейчас

\lp{20260604_194300166}

\noindent
же вышлем, если вам надо.

\section{23 августа 1914 г.}

\lp{20260604_194300166}

Дорогой мой!

Сегодня пришла ваша открытка Коле с поздравлением его ко дню рождения. Нам ее
переслали из Юдино. Сегодня у нас перебывала масса народу. Зяська, Саша Гедике,
Боря, Heimbach, Лев Эдуардович, Струве, супруги Гольденвейзер. Heimbach очень
тревожится за своего брата, о котором ничего ровно не знает. Я ее навещаю, но,
к сожалению, ничем не могу облегчить ее тяжелого положения. Лев Эдуардович
ходил заявляться, как это требуется для порядку,

\lp{20260604_194300166}

\noindent
но его в солдаты не взяли, и он, по-видимому, недоволен и рвется в бой.
Гольденвейзера освободили. Струве -- офицер запаса, и ему через несколько дней
вероятно придется вступить в ряды. Он бедный, так и не дождется, вероятно,
своих, которые застряли в Берлине, и уедет на службу.

Переживаем мы сейчас ужасно много, но не могу писать. Хоть бы знать, что эти
мои клочки доходят до вас. Я бы тогда писала гораздо больше. Коле, конечно, не
работается. Саша Гедике главным образом читает газеты и знает их чуть ли не
наизусть.

\lp{20260604_194241882}

До свидания, мой дорогой! С нетерпением жду известий. Ваша Анюта.

Это мое шестое письмо.
